% !TEX root = ../nc-quadric-surface-classification.tex
\subsection{階数三の四周期表示の生成元と関係}
\label{subsec:rank-three-relations}

\begin{proposition}
\label{prop:rank-three-relations}
定理~\ref{thm:normal-helix} の階数三の場合に、
矢を $E_i\to E_j$ の向きで書き、道の積を合成で定める。
隣接する二元の矢に加えて、$E_{2i}\to E_{2i+2}$ ごとに一つの追加矢を取る。
最小関係は以下の三種類である。
\begin{enumerate}
\item $E_{2i}\to E_{2i+2}$ に一つ。二つの隣接矢の空間を通る合成の核で、
その二因子テンソルの階数は二である。
\item $E_{2i}\to E_{2i+3}$ に二つ。
$X_{i+1}$ の極小評価列の核の次数一部分を、
$\Hom(\pi B(i),\pi B(i+1))$ と $\Hom(\pi B(i+1),\pi X_{i+1}(i+1))$
を経由する道として表す。
\item $E_{2i+1}\to E_{2i+4}$ に二つ。
$X_i$ の双対評価列の核の最初の非零次数部分を、
$\Hom(\pi X_i(i),\pi B(i+1))$ と $\Hom(\pi B(i+1),\pi B(i+2))$
を経由する道として表す。
\end{enumerate}
各距離三の関係は、距離二の関係を先に課した後の道の空間から
一意な二次元部分空間として得られる。
\end{proposition}

\begin{proof}
四周期表示の単純加群の極小分解
（系~\ref{cor:normal-cubic-defect} の証明）により、最小関係の個数は
偶数項間の距離二で一、距離三で各二であり、他にはない。
距離二の偶数項間では階数三の合成に第四の基底を加えてあるので、
唯一の関係の追加矢における係数は零である。
この関係が階数二になることは
注意~\ref{rem:rank-three-matrix} の行列分解の議論による。

次に $E_{2i}\to E_{2i+3}$ の道を考える。
距離二の関係を先に課すと、道の空間は
$\Hom(\pi B(i),\pi B(i+1))\otimes\Hom(\pi B(i+1),\pi X_{i+1}(i+1))$
と同定できる。実際、距離三には三つの隣接矢からなる道と
追加矢を一つ含む道があり、前者に距離二の関係を適用すれば
最初の二矢の積と追加矢を合わせて $B_1$ の基底が得られる。
従って合成の核を求めるには、$X_{i+1}$ の評価列
\eqref{eq:normal-ruling-evaluation} を次数一で取ればよい。
その核は $(p(X_{i+1}))_0$ であり、二次元である。

$E_{2i+1}\to E_{2i+4}$ についても、
後二つの隣接矢の積と追加矢を合わせると $B_1$ の基底が得られる。
ここで $B$ への内部 Hom を取った評価列
\[
 0\to\Hom_B(p(X_i),B)\to B^2
       \to \Hom_B(X_i,B)(1)\to0
\]
を使うと、合成の核は左端の次数一成分である。
これも二次元であり、主張した全ての距離三の関係を与える。
\end{proof}

比較先の四周期環を逆に構成するとき、
この表示で距離二の関係を階数二に正規化してから
距離三の各二次元部分空間を変化させられる。
ここから得る候補の全次数の許容条件は、
極小分解が実際に完全であることと、周期偶数側が
命題~\ref{prop:first-orbit-object} の軌道を与えることを含む。
これらは有限次元の関係の個数からは従わない。
